% !Mode:: "TeX:UTF-8"
\chapter*{总结与展望\markboth{总结与展望}{}}
\addcontentsline{toc}{chapter}{总结与展望}

\section*{工作总结}

本文研究了PyTorch深度学习框架模糊测试结果的管理与分析问题。
随着深度学习框架在实际系统中应用范围扩大，
模糊测试逐渐成为发现底层算子和参数边界缺陷的重要手段。
模糊测试会产生测试脚本结构、执行状态、覆盖贡献数据和崩溃日志等多类测试产物，
这些数据如果缺少统一组织，
后续的用例分析、故障分析和补充测试设计都需要较多人工整理。
鉴于上述问题，本文设计并实现了一套结合语义检索和大语言模型的
模糊测试数据管理与分析系统。
本文完成的工作如下：

1. 针对模糊测试产物类型较多、统一管理不便的问题，
实现了测试用例和测试产物的形式化定义，
采用结构化元数据和语义向量协同存储的方式建立测试用例之间的关联关系。
提出了基于双路向量化的相似用例召回方法，
将测试脚本的算子结构信息和崩溃日志文本分别编码为结构向量和日志向量，
并通过MySQL中保存的向量主键与Milvus中的向量集合完成跨库关联。
在此基础上，本文实现了测试用例的统一导入、列表查询、条件筛选、
详情查看和相似用例召回功能，
使原本分散的测试产物可以在同一条数据流程中被组织和检索，
为后续数据分析和故障分析提供数据基础。

2. 针对用例价值分布和补充测试方向难以直观分析的问题，
实现了用例效能评估任务的形式化定义，
提出了覆盖效能指数（Coverage Efficiency Index, CEI）模型。
该模型综合考虑测试用例的分支覆盖贡献和加权计算成本，
并使用动态百分位阈值筛选低收益用例。
物理执行环节使用\texttt{multiprocessing.spawn}进行进程隔离，
结合Gcov/LCOV采集PyTorch C++层的覆盖贡献数据，
降低单个异常用例影响主服务的风险。
在真实数据集上的实验结果表明，
在2000例模拟实验中，CEI筛选可以在覆盖率损失不足4\%的情况下
将总执行耗时缩减约60\%；
在200例物理对照实验中，物理命中行数保持率约为99.991\%，
总执行耗时缩减约33.2\%；
在1000例物理跑测中，物理命中行数仅下降0.005\%，
总执行耗时缩减约34.5\%。
在此基础上，系统进一步提供覆盖缺口分析、测试用例多样性分布、
用例效能排行、低效用例筛选和候选补盲测试脚本生成等功能，
用于辅助用户了解测试集覆盖情况和用例价值分布。

3. 针对崩溃用例中故障线索和候选边界信息不易整合分析的问题，
实现了基于DFS拓扑回溯的故障辅助定位方法，
对崩溃用例的算子调用链做反向遍历得到候选算子节点。
在相似样本召回的基础上，
实现了基于Daikon约束模板扫描的候选边界条件分析方法：
模块召回与目标崩溃用例算子拓扑相近的健康组和崩溃组样本，
将嵌套的张量参数展平成数值键值对，
再使用预设约束模板扫描候选不变量；
只有当某条约束在健康组样本中全部成立、在崩溃组样本中全部不成立时，
才被记为候选边界条件。
实验结果显示，该方法在五条已知算子约束上的召回率达到100\%，
并在一条包含\texttt{Conv1d}$\to$\texttt{GroupNorm}的真实崩溃用例上
给出了跨算子的隐式整除候选约束。
在此基础上，系统利用检索增强生成
（Retrieval-Augmented Generation, RAG）方法，
整理相似用例、运行信息和候选边界信息，
由LLM生成辅助诊断报告和候选补盲测试脚本，
为用户复查崩溃原因、理解参数边界和设计补充测试提供参考。

\section*{未来展望}

本文研究成果已经初步应用于PyTorch模糊测试结果的管理与分析任务，
但仍存在一些可改进之处和研究方向，具体有以下三点：

1. 在数据管理方面，测试数据规模仍有进一步丰富的空间。
本文主要基于2000例模拟数据、200例物理对照实验和1000例物理跑测做分析，
实验数据集中在有限的PyTorch版本和算子类型上。
未来的工作中，可以接入更多PyTorch版本、更多算子类型
和更长时间的Fuzzing任务数据，
以检验本文方法在不同运行环境下的适用性。

2. 在数据分析方面，CEI模型的算子权重仍有进一步优化的空间。
当前算子权重主要通过规则化方式设定，
不能完全等价于真实硬件执行成本。
未来的工作中，可以结合历史执行耗时、Profiling结果
以及资源占用情况，
对CEI模型中的权重和动态百分位阈值做更精细的调整。
同时，也可以将候选补盲测试脚本的生成结果与后续执行反馈结合，
形成更完整的覆盖缺口驱动的补充测试分析流程。

3. 在故障分析方面，候选边界条件分析和辅助诊断报告生成
仍依赖样本质量和人工复核。
如果相似健康样本和崩溃样本数量不足，
或参数分布中存在较多噪声，
候选约束可能出现漏检或误判；
RAG生成的辅助诊断报告也仍然需要源码级确认和真实复现。
未来的工作中，可以扩展约束模板库以支持更多跨算子约束类型，
并建立模型评分的基准，
完善候选补盲测试脚本的执行验证和回归测试流程，
从而提高故障分析结果的可解释性。
